\documentclass[11pt,a4paper]{article}

\usepackage[english, bulgarian]{babel}
\usepackage[T2A]{fontenc}
\usepackage{amssymb}
\usepackage{amsmath}
\usepackage{fancyhdr}

\ifdefined\showmargins
\PassOptionsToPackage{showframe}{geometry}
\fi
\usepackage[top=2cm, bottom=2cm, left=2cm, right=2cm]{geometry}
\usepackage[colorlinks=true, linkcolor=blue]{hyperref}
\usepackage{pgffor}

\pagestyle{fancy}
\fancyhf{}
\fancyhead[C]{СДРУЖЕНИЕ С НЕСТОПАНСКА ЦЕЛ „ПЕШАЧКО”}
\fancyhead[R]{\thepage}
\setlength{\headheight}{15pt}% to suppress fancyhdr warnings

\renewcommand{\labelenumi}{(\arabic{enumi})}

\newcounter{article}
\newcommand{\article}[2][]{
  \refstepcounter{article}% increment counter
  \noindent\textbf{Чл.~\arabic{article}.} #2
  \ifx\\#1\\\else\label{#1}\fi% set label if optional argument is given
}

\newcommand{\heading}[1]{
  \begin{center}
    \vspace{0.2cm}
    {\bf\large \MakeUppercase{#1}}
  \end{center}
  \addcontentsline{toc}{subsection}{#1}
}

\renewcommand{\thesection}{\Roman{section}}

\begin{document}

\begin{center}
  {\bf\Huge УСТАВ} \\
  на \\
  сдружение с нестопанска цел „ПЕШАЧКО”
\end{center}

\tableofcontents

\vspace{0.7cm}

\section{Наименование, цели, средства, предмет}

\heading{Статут}
\article{}
\begin{enumerate}
\item „ПЕШАЧКО” е сдружение (наричано по-долу „СДРУЖЕНИЕТО”) с нестопанска цел за
  осъществяване на дейност в обществена полза по смисъла на Закона за юридическите лица
  с нестопанска цел (ЗЮЛНЦ).
\item СДРУЖЕНИЕТО е юридическо лице отделно от членовете си и отговаря за задълженията
  си със своето имущество. Членовете на СДРУЖЕНИЕТО не отговарят лично за неговите
  задължения.
\end{enumerate}

\heading{Наименование}
\article{}
\begin{enumerate}
\item Наименованието на СДРУЖЕНИЕТО е: СДРУЖЕНИЕ С НЕСТОПАНСКА ЦЕЛ „ПЕШАЧКО”, което ще
  се изписва на латиница „PESHACHKO”.
\item Всяко писмено изявление от името на СДРУЖЕНИЕТО трябва да съдържа данните описани
  в {\bf Чл.~9} от ЗЮЛНЦ.
\end{enumerate}

\heading{Седалище и адрес}
\article{Седалището на СДРУЖЕНИЕТО е в Република България, гр. София, а адресът на
  управление -- гр.~София, п.к 1309 ж.к. „Света троица” бл.~379 вх.~Д. ет.~6 ап.~85.}

\heading{Цели}
\article[art:aim]{Основни цели на СДРУЖЕНИЕТО са:}
\begin{itemize}
\item Да популяризира ходенето като едно от най-важните условия за {\bf изграждане на
  здраве};
\item Да способства ходенето пеша да се утвърди като личен и обществен избор, който
  допринася за {\bf чист въздух и осъзнатост};
\item Да помогне за изграждането на {\bf навици у нашите деца} за активен начин на
  живот;
\item Да припомни, че разходката е начин за споделяне на време и {\bf създаване на
  общност}.
\end{itemize}

\heading{Средства за постигане на целите}
\article{Целите на СДРУЖЕНИЕТО ще се осъществяват със следните средства:}
\begin{itemize}
\item Личен пример от всеки член на СДРУЖЕНИЕТО;
\item Организиране, провеждане и участие в дискусии, семинари, лекции и курсове на теми,
  свързани с целите на СДРУЖЕНИЕТО;
\item Организиране на разходки и спортни събития като начин за приобщаване и
  популяризиране на принципите на здравословен начин на живот;
\item Управление на проекти и кандидатстване за тяхното финансиране чрез международни
  или български програми, съобразени с целите на СДРУЖЕНИЕТО;
\item Участие в инициативи и проекти на други неправителствени организации, граждански
  институции, университети и болници, съобразени с целите на СДРУЖЕНИЕТО;
\item Разработка на механизми за обмен на информация между различни институции и
  граждани от всички социални групи на обществото, във връзка с целите на СДРУЖЕНИЕТО;
\item Изграждане на устойчиви механизми за диалог с държавната администрация и
  насърчаване на пряко и конструктивно участие на гражданите в този диалог, съобразени с
  целите на СДРУЖЕНИЕТО;
\item Участие при планиране на изграждане на пешеходни пътни участъци и достъпен градски
  транспорт;
\item Иницииране и подпомагане на активни граждански дейности съобразени с целите на
  СДРУЖЕНИЕТО;
\item Популяризиране на целите и идеите на СДРУЖЕНИЕТО чрез брошури, електронни
  средства, публикации и кампании в медиите;
\item Създаване на прозрачна среда за комуникация и вземане на вътрешни за СДРУЖЕНИЕТО
  решения -- основополагащ елемент на конструктивно общуване с други сходни по идеи
  организации;
\item Използване на средства от членски внос и дарения, от допълнителна стопанска
  дейност, доброволен труд и други дейности, разрешени от закона;
\end{itemize}

\heading{Предмет на дейност}

\article{}
\begin{enumerate}
\item Предметът на дейност на СДРУЖЕНИЕТО е осъществяване на дейности в обществена
  полза, насочени към популяризиране на ходенето и активния начин на живот като средство
  за изграждане и поддържане на здраве, опазване на околната среда и изграждане на
  активна гражданска общност.

\item СДРУЖЕНИЕТО организира разходки, спортни и образователни събития, обучения,
  семинари, курсове, кампании и инициативи. Кандидатства по, участва в, и управлява
  проекти, както и сътрудничи на други организации и институции.

\item СДРУЖЕНИЕТО участва в процеси по планиране и обсъждане на изграждането на
  пешеходна инфраструктура и достъпен градски транспорт и също така разпространява
  информация и материали, свързани с целите си.

\item СДРУЖЕНИЕТО може да създава или да участва в създаването у нас и в чужбина на
  самостоятелни и съвместни специализирани звена, представителства, клонове,
  изследователски центрове, клубове, фондове и сдружения с нестопанска цел, по реда,
  определен от закона.

\item СДРУЖЕНИЕТО може да извършва и допълнителна стопанска дейност, свързана с предмета
  и целите си, при условията на закона.
\end{enumerate}

\heading{Стопанска дейност}
\article{}
\begin{enumerate}
\item СДРУЖЕНИЕТО може да извършва допълнителна стопанска дейност, свързана с предмета и
  целите му, при условията на закона. Като допълнителна стопанска дейност се включват:
  \begin{itemize}
  \item Организиране на платени обучения, семинари, курсове и други събития, свързани с
    целите на СДРУЖЕНИЕТО.
  \item Издателска, информационна и рекламна дейност, пряко свързани с дейността на
    СДРУЖЕНИЕТО.
  \item Предоставяне на консултантски услуги и обучения в областите, свързани с целите
    на СДРУЖЕНИЕТО;
  \item Управление и изпълнение на проекти, свързани с предмета и целите на СДРУЖЕНИЕТО;
  \item Отдаване под наем на движимо и недвижимо имущество;
  \item Други дейности, които са непосредствено свързани с целите на СДРУЖЕНИЕТО.
  \end{itemize}
\item Приходите от стопанската дейност се използват единствено за постигане на целите на
  СДРУЖЕНИЕТО.
\item СДРУЖЕНИЕТО не разпределя печалба.
\end{enumerate}

\heading{Източници на средства, имуществото}
\article{}
\begin{enumerate}
\item Средствата за осъществяване на дейността на СДРУЖЕНИЕТО се набавят от различни
  източници, включително:
  \begin{itemize}
  \item Членски внос на членовете на СДРУЖЕНИЕТО;
  \item Дарения, завещания, спонсорства и други доброволни вноски от физически и
    юридически лица;
  \item Приходи от проекти и програми, финансирани от национални или международни донори;
  \item Приходи от допълнителна стопанска дейност, извършвана в рамките на закона;
  \item Други източници, включително финансови инструменти, допустими по закон.
  \end{itemize}
\item Всички средства се използват изключително за постигане на целите на СДРУЖЕНИЕТО.
\end{enumerate}

\vspace{0.4cm}
\article{}

\begin{enumerate}
\item Имуществото на СДРУЖЕНИЕТО се формира от членския внос на членовете, имуществени
  вноски от членовете, дарения и завещания, спонсорство, субсидии и помощи, приходи от
  собствена допълнителна стопанска дейност, осъществявана с цел осигуряване
  самофинансирането на СДРУЖЕНИЕТО и постигане на неговите цели.
\item По решение на Общото събрание членовете на СДРУЖЕНИЕТО могат да правят целеви
  имуществени вноски в извънредни случаи, когато това се налага за провеждане на
  мероприятие във връзка с постигане на целите на СДРУЖЕНИЕТО.
\item Членовете на СДРУЖЕНИЕТО могат да му предоставят парични средства под формата на
  заем или временно и възмездно ползване на недвижими имоти и индивидуално определени
  движими вещи, като сключването на съответния договор се одобрява от Управителния съвет
  на СДРУЖЕНИЕТО.
\item Разходването на реализираните от СДРУЖЕНИЕТО парични средства и разпореждане с
  друго негово имущество се извършва за постигане на целите му по {\bf
    Чл.~\ref{art:aim}} от настоящия Устав.
\item СДРУЖЕНИЕТО може да разходва и безвъзмездно своето имущество, доколкото това е
  насочено към постигането на преследваните от него цели.
\end{enumerate}

\heading{Определяне на дейността}
\article{СДРУЖЕНИЕТО ще осъществява дейност в ОБЩЕСТВЕНА ПОЛЗА по смисъла на {\bf
    Чл.~2}, ал.~1 от ЗЮЛНЦ.}

\heading{Срок}
\article{СДРУЖЕНИЕТО не е ограничено със срок.}

\section{Членство}

\heading{Членство}
\article{}
\begin{enumerate}
\item Членове на СДРУЖЕНИЕТО могат да бъдат дееспособни физически лица и/или юридически
  лица, които заявят желание да спомогнат за постигане на целите и развитието на
  дейността на СДРУЖЕНИЕТО.
\item Членството в СДРУЖЕНИЕТО е доброволно, без ограничения по отношение на
  националност, етническа принадлежност, политически убеждения, раса, пол и религия.
\end{enumerate}

\heading{Придобиване на членство}
\article{}
\begin{enumerate}
\item Кандидатът подава заявление до Управителния съвет на СДРУЖЕНИЕТО, в което
  декларира, че е запознат и приема разпоредбите на настоящия Устав.
\item Кандидатите -- физически лица представят валиден електронен адрес (e-mail) за
  комуникация, както и техните имена. Кандидатите -- юридически лица представят
  заявление и преписи от документите си за регистрация и от решението на управителните
  си органи за членство в СДРУЖЕНИЕТО.
\item Управителният съвет разглежда заявлението за членство на кандидата в срок от 10
  работни дни от датата на неговото постъпване и с обикновено мнозинство (50\% + 1
  гласа) го приема или отхвърля. Отказът на Управителния съвет може да се обжалва пред
  Общото събрание. Членството в СДРУЖЕНИЕТО се придобива от датата на решението на
  Управителния съвет или от решението на Общото събрание.
\end{enumerate}

\heading{Права и задължения на членовете}
\article{Членовете на СДРУЖЕНИЕТО имат следните права:}
\begin{itemize}
\item да участват в управлението на СДРУЖЕНИЕТО, като избират Управителен съвет от 3
  лица;
\item да бъдат информирани за неговата дейност;
\item да ползват имуществото на СДРУЖЕНИЕТО за постигане целите му;
\item да ползват резултатите от дейността на СДРУЖЕНИЕТО, както и събраната информация,
  съгласно разпоредбите на този Устав.
\end{itemize}

\article{Членовете на СДРУЖЕНИЕТО са длъжни:}
\begin{itemize}
\item да спазват разпоредбите на този Устав и да изпълняват решенията на ръководните
  органи на СДРУЖЕНИЕТО;
\item да съдействат за изпълнението на неговите цели;
\item да издигат авторитета на СДРУЖЕНИЕТО;
\item да не извършват действия или бездействия, които противоречат на целите му, или да
  го злепоставят;
\item да внасят в срок предвидените в настоящия Устав членски вноски.
\end{itemize}

\article{Членските права и задължения, с изключение на имуществените, са непрехвърлими и
  не преминават върху други лица при смърт, съответно при прекратяване.}

\vspace{0.4cm}
\article{Членовете на СДРУЖЕНИЕТО имат право да овластят трето лице да упражнява техните
  права по конкретен повод и да изпълнява задълженията им, което се извършва писмено и
  произвежда действие само след писмено уведомяване на Управителния съвет. В този случай
  те носят отговорност за неизпълнение или изпълнение на техните задължения от страна на
  овластеното лице.}

\heading{Прекратяване на членството}
\article{Членството в СДРУЖЕНИЕТО се прекратява:}
\begin{itemize}
\item с едностранно писмено волеизявление, отправено до Управителния съвет на
  СДРУЖЕНИЕТО;
\item със смърт или поставяне под запрещение, респективно прекратяване на юридическата
  личност на член на СДРУЖЕНИЕТО;
\item с изключване;
\item с прекратяване на СДРУЖЕНИЕТО.
\end{itemize}

\article{}
\begin{enumerate}
\item Член на СДРУЖЕНИЕТО може да бъде изключен с решение на Управителния съвет по
  негова инициатива или след постъпване на писмено предложение от 1/3 от членовете на
  СДРУЖЕНИЕТО, когато:
  \begin{itemize}
  \item нарушава предвидените в този Устав задължения;
  \item извършва други действия, които правят по-нататъшното му членство в СДРУЖЕНИЕТО
    несъвместимо.
  \end{itemize}
\item При по-маловажни случаи на нарушения, Управителният съвет определя с решение срок
  за преустановяване на нарушението и за отстраняването на неговите последици.
\end{enumerate}

\article{При прекратяване на членството СДРУЖЕНИЕТО не дължи връщане на направените
  членски вноски.}

\heading{Членски внос}
\article{}
\begin{enumerate}
\item Членският внос се определя в размер на €25 за физически лица и половин минимална
  работна заплата за юридически лица.
\item Членският внос се внася ежегодно преди провеждането на годишното общо събрание.
\end{enumerate}

\section{Управление}

\heading{Органи на управление}
\article{Върховен орган на СДРУЖЕНИЕТО е Общото събрание, негов орган е Управителният
  съвет.}

\heading{Общо събрание}
\article{}
\begin{enumerate}
\item Общото събрание се състои от всички членове на СДРУЖЕНИЕТО.
\item Членовете на СДРУЖЕНИЕТО. участват в Общото събрание лично или чрез упълномощен
  представител.
\item Членовете – физически лица участват директно, а членовете – юридически лица се
  представляват в Общото събрание от законните им представители или от изрично
  упълномощено лице.
\item Пълномощник на юридическо лице или на физическо лице може да бъде само физическо
  лице.
\item Упълномощаването се извършва с изрично писмено пълномощно за конкретно Общо
  събрание, с посочени решения по дневния ред.
\item Пълномощниците нямат право да представляват повече от един член едновременно и да
  преупълномощават трети лица с правата си.
\end{enumerate}

\heading{Правомощия на Общото събрание}
\article[art:general-assembly-rights]{Общото събрание:}
\begin{itemize}
\item Изменя и допълва Устава.
\item Взема решение за преобразуване или прекратяване на СДРУЖЕНИЕТО.
\item Взема решение за участие в други организации.
\item Взема решение за разходване на парични средства в размер над €10'000, както и
  договори за заем/кредит от СДРУЖЕНИЕТО като заемополучател/кредитополучател на
  стойност над €10'000.
\item Взема решение за сделки с недвижими имоти.
\item Взема решение за за безвъзмездно разходване на всякакъв вид имуществото на
  СДРУЖЕНИЕТО, когато е в полза на:
  \begin{itemize}
  \item лица от състава на другите органи на СДРУЖЕНИЕТО и техните съпрузи, роднините им
    по права линия - без ограничение, по съребрена линия - до четвърта степен, или по
    сватовство - до втора степен включително;
  \item лица, били в състава на управителните му органи до 2 години преди датата на
    вземане на решението;
  \item юридически лица, в които посочените лица в горните две точки са управители или могат да
    наложат или възпрепятстват вземането на решения;
  \item юридически лица, финансирали СДРУЖЕНИЕТО до 3 години преди датата на вземане на решение;
  \item политически партии, в ръководните и контролните органи на които участват членове
    на ръководни и контролни органи на СДРУЖЕНИЕТО.
  \end{itemize}
\end{itemize}

\article{Общото събрание:}
\begin{itemize}
\item Избира и освобождава членовете на Управителния съвет.
\item Приема основните насоки и програми за дейността на СДРУЖЕНИЕТО.
\item Приема бюджета на СДРУЖЕНИЕТО.
\item Произнася се с решение по жалби на изключени и кандидат-членове.
\item Приема отчета за дейността на Управителния съвет.
\item Отменя решенията на другите органи на СДРУЖЕНИЕТО, които противоречат на закона,
  Устава или други вътрешни актове, регламентиращи дейността на СДРУЖЕНИЕТО.
\item Взема решения по всички други въпроси, поставени в неговата компетентност от
  закона или настоящия Устав.
\end{itemize}

\heading{Свикване}
\article{Общото събрание се свиква на редовно заседание от Управителния съвет на
  СДРУЖЕНИЕТО, веднъж годишно.}

\vspace{0.4cm}
\article{}
\begin{enumerate}
\item Общото събрание се свиква с писмени покани или чрез покана, изпратена в електронна
  форма на всеки член на Общото събрание, изпратена на електронен адрес, посочен от него
  за контакт, най-късно 7 дни преди датата на заседанието.
\item Поканата трябва да съдържа дневния ред, датата, часа, мястото и начина за
  провеждането на общото събрание и по чия инициатива то се свиква.
\end{enumerate}

\vspace{0.4cm}
\article{Извънредно заседание на Общо събрание се свиква по инициатива на Управителния
  съвет. Една трета от членовете на СДРУЖЕНИЕТО имат право да поискат Управителния съвет
  да свика Общо събрание по конкретен въпрос и ако той не отправи покана в двуседмичен
  срок от датата на поискването, събранието се свиква от съда по седалището на
  СДРУЖЕНИЕТО.}

\vspace{0.4cm}
\article{По изключение Общото събрание може да се проведе и онлайн чрез конферентна
  интернет връзка или по друг подходящ начин, който гарантира установяване на
  самоличността на лицата, участващи в Общото събрание, както и възможност те да
  участват в обсъждането и вземането на решения.}

\heading{Кворум}
\article{}
\begin{enumerate}
\item Общото събрание е законно, ако присъстват повече от половината от членовете на
  СДРУЖЕНИЕТО. При липса на кворум, присъстващите могат да решат с обикновено мнозинство
  отлагане за друга дата или прилагане на процедурата описана в {\bf Чл.~27} от ЗЮЛНЦ.
  Общо събрание не може да бъде отлагано повече от веднъж.
\item Кворумът се установява от председателстващия събранието и присъстващите се
  отразяват в списък, който се прилага към протокола на заседанието.
\end{enumerate}

\heading{Гласуване}
\article{При гласуване всеки член на СДРУЖЕНИЕТО има право само на един глас.}

\heading{Решения на общото събрание}
\article{Решенията на Общото събрание се вземат с обикновено мнозинство от
  присъстващите. Решенията по {\bf Чл.~\ref{art:general-assembly-rights}}, се вземат с
  квалифицирано мнозинство (2/3 от присъстващите).}

\vspace{0.4cm}
\article{Общото събрание не може да взема решения, които не са включени в предварително
  обявения дневен ред на събранието.}

\heading{Протокол}
\article{}
\begin{enumerate}
\item За всяко заседание на Общото събрание се води протокол, в който се отразяват
  направените разисквания, предложения и взетите решения. Протоколите се удостоверяват с
  подписите на председателя и на секретаря-протоколчик и се съхраняват в специални книги
  (или в електронен формат).
\item Всеки член, присъствал на Общото събрание има право да следи за точното отразяване
  на заседанието и взетите решения в протокола, както и да получи достъп до протоколи от
  минали заседания.
\end{enumerate}

\heading{Управителен съвет}
\article{Управителният съвет се състои от 3 лица -- членове на СДРУЖЕНИЕТО.}

\heading{Мандат}
\article{Управителният съвет се избира за срок от 3 години, като неговите членове могат
  да бъдат преизбирани неограничено.}

\heading{Правомощия на Управителния съвет}
\article{Управителният съвет:}
\begin{itemize}
\item Представлява СДРУЖЕНИЕТО и определя обема на представителната власт, на
  председателя и секретаря.
\item Избира из между членовете си председател и секретар.
\item Осигурява изпълнението на решенията на Общото събрание.
\item Взема решение за откриване и закриване на клонове и назначава управител/и.
\item Подготвя и внася в Общото събрание проект за бюджет.
\item Подготвя и внася в Общото събрание отчет за дейността на СДРУЖЕНИЕТО.
\item Определя реда и организира извършването на дейността на СДРУЖЕНИЕТО и носи
  отговорност за това.
\item Определя адреса на СДРУЖЕНИЕТО.
\item Приема правила за работата си и пряко ръководи стопанската дейност на СДРУЖЕНИЕТО.
\item Взема решения по всички въпроси, освен тези които са от компетентността на Общото
  събрание на СДРУЖЕНИЕТО.
\item Приема и изключва членове на СДРУЖЕНИЕТО.
\end{itemize}

\heading{Заседания}
\article{Редовните заседания на Управителния съвет се свикват от председателя по негова
  инициатива веднъж месечно, както и по искане на всеки един от неговите членове. Ако
  председателят не свика заседание в седмичен срок от поискването, такова може да се
  свика от всеки един от членовете на Управителния съвет.}

\vspace{0.4cm}
\article{}
\begin{enumerate}
\item Заседанието е законно, ако на него присъстват повече от половината от членовете на
  Управителния съвет.
\item Законно решение може да бъде взето и без да се провежда заседание, ако протоколът
  за това бъде подписан без забележки и възражения за това от всички членове на
  Управителния съвет.
\item Заседанията се ръководят от председателя на Управителния съвет, а в негово
  отсъствие от посочен от него член.
\end{enumerate}

\heading{Решения на Управителния съвет}
\article{Управителният съвет взема решенията си с обикновено мнозинство от присъстващите
  членове.}

\heading{Председател на управителния съвет}
\article{}
\begin{enumerate}
\item Управителният съвет избира от своя състав председател. Решението на Управителния
  съвет се взема с мнозинство 2/3 от всички членове. Председателският мандат е 3 (три)
  години, освен ако по обективни причини същият не бъде прекратен от Управителния съвет,
  респективно Председателят не бъде освободен като член на Управителния съвет.
\item Председателят на СДРУЖЕНИЕТО:
  \begin{itemize}
  \item Представлява СДРУЖЕНИЕТО, в съответствие с решенията на Управителния съвет;
  \item Взаимодейства с органите на държавна и изпълнителна власт, български и
    чуждестранни институции и структури;
  \item Допринася за утвърждаване на високия авторитет и публичния облик на СДРУЖЕНИЕТО;
  \item Организира дейността на СДРУЖЕНИЕТО и осъществява оперативното му ръководство;
  \item Организира изпълнението на решенията на Управителния съвет;
  \item Сключва трудовите договори със служителите на СДРУЖЕНИЕТО и определя трудовите
    им възнаграждения, освен с тези, които се назначават от Управителния съвет;
  \item Възлага на трети лица извършването на определени действия и/или изготвянето на
    определени документи в интерес или за целите на СДРУЖЕНИЕТО;
  \item Докладва незабавно на Управителния съвет за съществени обстоятелства, касаещи
    дейността на СДРУЖЕНИЕТО.
  \end{itemize}
\end{enumerate}

\heading{Контрол}
\article{Контрол се осъществява от Общото събрание на СДРУЖЕНИЕТО.}

\heading{Финансов отчет}

\article{}
\begin{enumerate}
\item Управителният съвет е длъжен периодично да изготвя предвидената в Закона за
  счетоводството отчетна информация за дейността на СДРУЖЕНИЕТО при спазване на
  принципите за откритост, достоверност и своевременност.
\item Счетоводството на СДРУЖЕНИЕТО трябва да бъде достъпно за всички негови членове.
\end{enumerate}

\vspace{0.4cm}
\article{}
\begin{enumerate}
\item Всяка година Управителният съвет осигурява съставянето на годишен финансов отчет и
  годишен доклад за дейността за изтеклата календарна година.
\item Годишният финансов отчет на СДРУЖЕНИЕТО подлежи на независим финансов одит, при
  спазване на реда и ако са налице условията, предвидени в Закона за счетоводството.
\end{enumerate}

\heading{Клонове}
\article{С решение на Управителния съвет на СДРУЖЕНИЕТО могат да се откриват и закриват
  клонове извън населеното място, в което се намира седалището на СДРУЖЕНИЕТО.}

\vspace{0.4cm}
\article{Клоновете не са юридически лица, ръководят се от управител, назначен от
  Управителния съвет и извършват дейностите, определени с решението на Управителния
  съвет. Със същото решение се определят и ограниченията в правомощията им и
  представителната власт на управителя на клона.}

\vspace{0.4cm}
\article{Клоновете водят книги за дейността си и веднъж на месец управителят на клона
  представя пред Управителният съвет на СДРУЖЕНИЕТО отчет за дейността на клона и
  разходваните средства.}

\vspace{0.4cm}
\article{При откриване на клон, Управителният съвет трябва да опише в протокол,
  седалището и адреса на клона, неговия управител и ограниченията на неговите
  правомощия. На заявление подлежат и промените в горепосочените обстоятелства.
  Решението за откриване на клон трябва да е обявено публично и протоколът да е достъпен
  за всички членове на СДРУЖЕНИЕТО.}

\section{Прекратяване и ликвидация}

\heading{Прекратяване}
\article{СДРУЖЕНИЕТО се прекратява с:}
\begin{enumerate}
\item с решение на Общото събрание;
\item с решение на Окръжния съд по седалището в случаите описани в {\bf Чл.~13}, ал.~1, т.~3
  от ЗЮЛНЦ.
\end{enumerate}

\heading{Ликвидация}
\article{}
\begin{enumerate}
\item В случай на прекратяване на СДРУЖЕНИЕТО се извършва ликвидация от Управителния
  съвет или от определено от него лице.
\item Ликвидаторът е длъжен по възможност да удовлетвори кредиторите на СДРУЖЕНИЕТО от
  наличните парични средства, а ако това е невъзможно -- чрез осребряване първо на
  движимото, а след това на недвижимото имущество на юридическото лице с нестопанска
  цел.
\item Имущество не може да се прехвърля по какъвто и да е начин на:
  \begin{enumerate}
  \item учредителите и настоящите и бившите членове;
  \item лицата, били в състава на органите му и служителите му;
  \item ликвидаторите освен дължимото възнаграждение;
  \item съпрузите на лицата по точки (а), (б) и (в);
  \item роднините на лицата по точки (а), (б) и (в) по права линия - без ограничение, по
    съребрена линия - до четвърта степен, или по сватовство - до втора степен
    включително;
  \item юридическите лица, в които лицата по точки (а), (б), (в), (г) и (д) са
    управители или могат да наложат или възпрепятстват вземането на решения.
  \end{enumerate}
\item Имуществото, останало след удовлетворението на кредиторите, се предоставя на
  юридическо лице със същата или близка нестопанска цел, избрано от Общото събрание с
  решението за откриване на производство по ликвидация.
\end{enumerate}

\section{Заключителни разпоредби}
\article{}
\begin{enumerate}
\item Промени в настоящия Устав могат да бъдат извършвани по реда, предвиден в него и в
  ЗЮЛНЦ.
\item За всички неуредени с настоящия Устав въпроси се прилага действащото българско
  законодателство.
\item В случай, че някоя от разпоредбите на настоящия Устав е или стане недействителна
  или неприложима, това не влияе на действителността на останалите клаузи.
  Недействителната или неприложима клауза следва да бъде заместена от императивна
  законова разпоредба или от такава действителна или приложима клауза, с която ще бъдат
  постигнати оптимално целта и намеренията на учредителите.
\item Този устав е приет на учредително събрание за създаване на сдружение с нестопанска
  цел „ПЕШАЧКО”, проведено в гр.~София, п.к 1309 ж.к. „Света троица” бл.~379 вх.~Д.
  ет.~6 ап.~85.
\item Списъкът на членовете на СДРУЖЕНИЕТО, подписали устава се намира в
  Раздел~\nameref{sec:members} и се счита за неразделна част от този Устав.
\end{enumerate}

\newpage

\section{СПИСЪК НА ЧЛЕНОВЕТЕ НА СДРУЖЕНИЕ С НЕСТОПАНСКА ЦЕЛ „ПЕШАЧКО”}\label{sec:members}

\vspace{1cm} { \renewcommand{\labelenumi}{\arabic{enumi}.}
\begin{enumerate}
  \foreach \i in {1,...,11}{
  \item \makebox[10cm]{\dotfill}, ЕГН: \makebox[4.8cm]{\dotfill} \\[0.4cm]
  }
\end{enumerate}
}

\end{document}
